% Part: lambda-calculus
% Chapter: syntax
% Section: substitution

\documentclass[../../../include/open-logic-section]{subfiles}

\begin{document}

\olfileid{lam}{syn}{sub}
\olsection{आदेशन}

\begin{explain}
मुक्त चरे हे परिसरातील चरांचे संदर्भ असतात. म्हणून मुक्त
चराच्या जागी ठराविक मूल्य वापरणे अर्थपूर्ण ठरते. उदाहरणार्थ,
$\lambd[x][f x]$ मधील $f$ च्या जागी ओळख फलनासारखे
$\lambd[y][y]$ हे विशिष्ट पद ठेवायचे असू शकते. त्यातून
$\lambd[x][(\lambd[y][y]) x]$ मिळते. मुक्त चरांच्या जागी
लॅम्डा पदे ठेवण्याच्या प्रक्रियेला आदेशन म्हणतात.
\end{explain}

\begin{defn}[आदेशन] \ollabel{defn:substitution}
  पद $M$ मधील चर $x$ च्या जागी पद $N$ चे
  \emph{आदेशन}, $\Subst{M}{N}{x}$, आगमनाने पुढीलप्रमाणे
  परिभाषित करतात:
  \begin{enumerate}
    \item $\Subst{x}{N}{x} = N$. \ollabel{defn:substitution-1}
    \item $\Subst{y}{N}{x}  = y$, जर $x \neq y$. \ollabel{defn:substitution-2}
    \item $\Subst{PQ}{N}{x}  = (\Subst{P}{N}{x}) (\Subst{Q}{N}{x})$.
      \ollabel{def:substitution-3}
    \item $\Subst{(\lambd[y][P])}{N}{x} = \lambd[y][\Subst{P}{N}{x}]$,
      जर $x \neq y$ आणि $y \notin \FV{N}$; अन्यथा अपरिभाषित.
      \ollabel{defn:substitution-4}
  \end{enumerate}
\end{defn}

\begin{explain}
\olref{defn:substitution}\olref{defn:substitution-4} मध्ये
$x \neq y$ ही अट घालण्याचे कारण असे की~$M$ मधील
चर~$x$ च्या \emph{बद्ध} आस्थितींच्या जागी~$N$ ठेवायचे नाही.
उदाहरणार्थ, $\Subst{\lambd[x][x]}{y}{x}$ विचारात घेताना
$\lambd[x][y]$ मिळू नये; बद्ध आस्थिती बदलली नाही तर
$\lambd[x][x]$ हेच पद राहते. मात्र वरच्या औपचारिक
नियमांनुसार या प्रसंगी आदेशन अपरिभाषित आहे.

$\lambd[y][P]$ मध्ये~$x$ च्या जागी $N$ चे आदेशन करताना
$y \notin \FV{N}$ हीही अट घालतात. उदाहरणार्थ,
$\lambd[y][x]$ मध्ये $x$ च्या जागी $y$ चे आदेशन,
म्हणजे $\Subst{\lambd[y][x]}{y}{x}$, करता येत नाही:
त्यातून $\lambd[y][y]$ मिळेल. हे पद युक्तिवाद स्वीकारून
तोच परत करणारे फलन दर्शवते. पण $\lambd[y][x]$ हे पद
नेहमी~$x$ (किंवा $x$ ज्याचा संदर्भ देते ते) परत करणारे
फलन दर्शवते. म्हणून अपेक्षित परिणाम म्हणजे युक्तिवाद
स्वीकारून तो बाजूला ठेवणारे आणि परिसरातील चर~$y$
परत करणारे फलन. हे नीट साधण्यासाठी आधी बद्ध चर~$y$
चे नाव ``बदलावे'' लागेल.
\end{explain}

\begin{prob}
  पुढील आदेशनांचे परिणाम काय आहेत?
  \begin{enumerate}
  \item $\Subst{\lambd[y][x(\lambd[w][vwx])]}{(uv)}{x}$
  \item $\Subst{\lambd[y][x(\lambd[x][x])]}{(\lambd[y][xy])}{x}$
  \item $\Subst{y(\lambd[v][xv])}{(\lambd[y][vy])}{x}$
  \end{enumerate}
\end{prob}

\begin{thm} \ollabel{thm:notinfv}
  $x \notin \FV{M}$ असेल, तर डावी बाजू परिभाषित
  असल्यास $\FV{\Subst{M}{N}{x}} = \FV{M}$.
\end{thm}

\begin{proof}
  $M$ च्या रचनेवर विगमन करू.
  \begin{enumerate}
  \item $M$ हे चर आहे: सरावासाठी.
  \item $M$ हे $(PQ)$ या रूपाचे आहे: सरावासाठी.
  \item $M$ हे $\lambd[y][P]$ या रूपाचे आहे. तसेच
    $\Subst{\lambd[y][P]}{N}{x}$ परिभाषित असल्यामुळे
    ते $\lambd[y][\Subst{P}{N}{x}]$ असते. म्हणून
    $\Subst{P}{N}{x}$ परिभाषित असते; शिवाय $x \neq y$
    आणि $x \notin \FV{P}$. त्यामुळे:
    \begin{multline*}
      \FV{\Subst{\lambd[y][P]}{N}{x}} = \\
    \begin{aligned}
      & = \FV{\lambd[y][\Subst{P}{N}{x}]} &&
      \text{\olref{defn:substitution-4} नुसार}\\
      & = \FV{\Subst{P}{N}{x}} \setminus \{y\} &&
      \text{\olref[fv]{def:fv}\olref[fv]{def:fv2} नुसार}\\
      & = \FV{P} \setminus \{y\} && \text{विगमन गृहीतकानुसार} \\
      & = \FV{\lambd[y][P]} && \text{\olref[fv]{def:fv}\olref[fv]{def:fv2} नुसार}
    \end{aligned}
    \end{multline*}
  \end{enumerate}
\end{proof}

\begin{prob}
  \olref[lam][syn][sub]{thm:notinfv} ची सिद्धता पूर्ण करा.
\end{prob}

\begin{thm} \ollabel{thm:infv}
  $x \in \FV{M}$ असेल, तर डावी बाजू परिभाषित
  असल्यास $\FV{\Subst{M}{N}{x}} = (\FV{M} \setminus
  \{x\}) \cup \FV{N}$.
\end{thm}

\begin{proof}
  $M$ च्या रचनेवर विगमन करू.
  \begin{enumerate}
  \item $M$ हे चर आहे: सरावासाठी.
  \item $M$ हे $PQ$ या रूपाचे आहे: $\Subst{(PQ)}{N}{x}$
    परिभाषित असल्यामुळे ते
    $(\Subst{P}{N}{x})(\Subst{Q}{N}{x})$ असते आणि दोन्ही
    आदेशन परिभाषित असतात. शिवाय $x \in \FV{PQ}$
    असल्यामुळे $x \in \FV{P}$ किंवा $x \in \FV{Q}$,
    किंवा दोन्ही, असते. उरलेला भाग सरावासाठी आहे.
  \item $M$ हे $\lambd[y][P]$ या रूपाचे आहे.
    $\Subst{\lambd[y][P]}{N}{x}$ परिभाषित असल्यामुळे
    ते $\lambd[y][\Subst{P}{N}{x}]$ असते;
    $\Subst{P}{N}{x}$ परिभाषित असते, $x \neq y$ आणि
    $y \notin \FV{N}$. तसेच $x \in
    \FV{\lambd[y][P]}$ असल्यामुळे $x \in \FV{P}$.
    आता:
    \begin{multline*}
      \FV{\Subst{(\lambd[y][P])}{N}{x}} = \\
      \begin{aligned}
      & = \FV{\lambd[y][\Subst{P}{N}{x}]} \\
      & = \FV{\Subst{P}{N}{x}} \setminus \{y\} \\
      & = ((\FV{P} \setminus \{x\}) \cup \FV{N}) \setminus \{y\}
       && \text{विगमन गृहीतकानुसार} \\
      & = (\FV{P} \setminus \{x, y\}) \cup \FV{N}
       && y \notin \FV{N} \\
      & = (\FV{\lambd[y][P]} \setminus \{x\}) \cup \FV{N}
      \end{aligned}
      \end{multline*}
  \end{enumerate}
\end{proof}

\begin{prob}
  \olref[lam][syn][sub]{thm:infv} ची सिद्धता पूर्ण करा.
\end{prob}

\begin{thm}\ollabel{thm:clr}
  $x \notin \FV{N}$ आणि आदेशन परिभाषित
  असेल, तर $x \notin \FV{\Subst{M}{N}{x}}$.
\end{thm}

\begin{proof}
  सरावासाठी.
\end{proof}

\begin{prob}
  \olref[lam][syn][sub]{thm:clr} सिद्ध करा.
\end{prob}


\begin{thm}\ollabel{thm:inv}
  $\Subst{M}{y}{x}$ परिभाषित असेल आणि $y \notin \FV{M}$,
  तर $\Subst{\Subst{M}{y}{x}}{x}{y} = M$.
\end{thm}

\begin{proof}
  $M$ च्या रचनेवर विगमन करू.
  \begin{enumerate}
  \item $M$ हे चर $z$ आहे: सरावासाठी.
  \item $M$ हे $(PQ)$ या रूपाचे आहे. मग:
    \begin{align*}
      \Subst{\Subst{(PQ)}{y}{x}}{x}{y}
      &=\Subst{((\Subst{P}{y}{x})(\Subst{Q}{y}{x}))}{x}{y} \\
      &= (\Subst{\Subst{P}{y}{x}}{x}{y})(\Subst{\Subst{Q}{y}{x}}{x}{y}) \\
      &= (PQ) \text{ विगमन गृहीतकानुसार}
    \end{align*}
  \item $M$ हे $\lambd[z][N]$ या रूपाचे आहे.
    $\Subst{\lambd[z][N]}{y}{x}$ परिभाषित असल्यामुळे
    $z \neq y$ हे माहीत आहे. म्हणून:
    \begin{align*}
      \Subst{\Subst{(\lambd[z][N])}{y}{x}}{x}{y}\\
      & = \Subst{(\lambd[z][\Subst{N}{y}{x}])}{x}{y} \\
      & = \lambd[z][\Subst{\Subst{N}{y}{x}}{x}{y}] \\
      &= \lambd[z][N] \text{ विगमन गृहीतकानुसार}
    \end{align*}
  \end{enumerate}
\end{proof}

\begin{prob}
  \olref[lam][syn][sub]{thm:inv} ची सिद्धता पूर्ण करा.
\end{prob}

\end{document}
